% Stage 8 restored fragment — 2.5.2 Near, intermediate, and far regimes (archive lines 6551-6826)
\subsection{Near, intermediate, and far regimes}
\label{sec:ch252}

The far-field approximation replaces distance-dependent amplitudes by $e^{-jkr}/r$ and angular factors \cite{jackson1999}. Range calibrations must sit in the Fraunhofer zone for this step \cite{hansen1988}. Maxwell's retarded export to $\Gamma_R$ remains the governing constraint \cite{maxwell1873}.
Multipole expansions order far-zone fields by angular complexity: the $\ell=1$ dipole term
dominates electrically small feeds, whereas reflector campaigns require many $\ell$ before
sidelobe floors stabilize \cite{stratton1941}. The far-field approximation replaces
distance-dependent amplitudes by $e^{-\jj k_0 r}/r$ and angular factors; range calibrations
must sit in the Fraunhofer zone for that step to be valid \cite{jackson1999}. A pattern
meter on $\Gamma_R$ with $k_0 R\ll 1$ still reads $\mathcal{S}_\omega[\Jimp]$; only when
$k_0 R\gg 1$ and $R\gg a$ does the same probe approximate $\mathcal{R}_\omega[\Jimp]$ through the
far-zone reduction developed below \cite{harrington2001}.

The same source produces all three regions simultaneously; the distinction arises from how
different terms of the radiation integral scale with distance \cite{stratton1941}. Near the
feed, reactive \(1/r^{3}\) and \(1/r^{2}\) terms dominate shell integrals; at intermediate
distances those terms compete with the leading \(e^{-\jj k_0 r}/r\) radiation; only on
sufficiently large spheres does the transverse \(1/r\) term alone support pattern and gain
definitions \cite{harrington2001}.
Let \(\Omega_s=\operatorname{supp}\Jimp\) be compact with characteristic scale \(a\). For
observation radius \(r=|\mathbf{r}|\) from a reference point in \(\Omega_s\), define
\begin{equation}
\label{eq:ch2.distance-regime}
\begin{aligned}
\text{near:}\ & k_0 r\lesssim 1,\\
\text{intermediate:}\ & 1\lesssim k_0 r\lesssim (k_0 a)^{-1},\\
\text{far:}\ & k_0 r\gg k_0 a,\quad r\gg a,
\end{aligned}
\end{equation}
In Eq.~\eqref{eq:ch2.distance-regime},
\(r=|\mathbf{r}|\) is the observation radius (m) from a reference point in
\(\Omega_s=\operatorname{supp}\Jimp\);
\(a\) is the source characteristic scale (m);
\(k_0=\omega\sqrt{\varepsilon_0\mu_0}\) is the vacuum wavenumber (rad\,m\(^{-1}\)); and
the inequalities label reactive near, transition, and radiative far bands before the
constants \(C_{\mathrm{n}},C_{\mathrm{i}},C_{\mathrm{f}}\) are fixed in
Eq.~\eqref{eq:ch2.distance-regime-definitions}.
with implicit constants fixed below \cite{stratton1941}. Equations~\eqref{eq:ch2.distance-regime}
are scale labels, not sharp cutoffs: reactive and radiative terms coexist in the near zone, and
intermediate zones require explicit asymptotic bookkeeping before source integrals reduce to
transverse \(1/r\) radiation \cite{harrington2001}. Pattern and gain definitions used in antenna
theory attach only after the far-zone reduction in
Eq.~\eqref{eq:ch2.far-zone-transverse-radiation} \cite{harrington2001}.
Field-magnitude scalings, far-zone reduction of
Eqs.~\eqref{eq:ch2.outgoing-electric-volume-integral}--\eqref{eq:ch2.combined-impressed-source-integral},
Fresnel--Fraunhofer transition scales, and power-flux convergence on large spheres are
developed below \cite{stratton1941}. Sommerfeld asymptotics
Eqs.~\eqref{eq:ch2.sommerfeld-dyadic-asymptotic} and~\eqref{eq:ch2.far-zone-transverse-radiation}
are presupposed; regime indices
Eqs.~\eqref{eq:ch1.regime-index}--\eqref{eq:ch1.radiative-regime-condition} are invoked by
reference \cite{stratton1941}.
Let \(\Omega_s=\operatorname{supp}\Jimp\) be compact with characteristic radius
\(a=\max_{\mathbf{r},\mathbf{r}'\in\Omega_s}|\mathbf{r}-\mathbf{r}'|/2\), vacuum
wavenumber \(k_0=\omega\sqrt{\varepsilon_0\mu_0}\), wavelength \(\lambda_0=2\pi/k_0\),
and electrical size \(k_0 a\) \cite{stratton1941}. Observation points
\(\mathbf{r}\) satisfy \(|\mathbf{r}|=r\notin\Omega_s\); reference origin is chosen in
\(\Omega_s\). Outgoing kernels from the source integrals above and dyadic structure from
Eq.~\eqref{eq:ch2.free-space-dyadic-green} are presupposed without re-tabulation
\cite{stratton1941}. Constants \(C_1,C_2,\ldots\) below depend on \(k_0,a\) and source
geometry but not on \(r\) in the asymptotic limits stated \cite{harrington2001}.
The locked regime partition is
\begin{equation}
\label{eq:ch2.distance-regime-definitions}
\begin{aligned}
\text{\emph{near (reactive) zone:}}\ & k_0 r\le C_{\mathrm{n}},\\[2pt]
\text{\emph{intermediate (transition) zone:}}\ &
C_{\mathrm{n}}<k_0 r\le C_{\mathrm{i}}(k_0 a),\\[2pt]
\text{\emph{far (radiative) zone:}}\ &
k_0 r> C_{\mathrm{i}}(k_0 a),\quad r> C_{\mathrm{f}}a^{2}/\lambda_0,
\end{aligned}
\end{equation}
with declared positive constants \(C_{\mathrm{n}}=\mathcal{O}(1)\),
\(C_{\mathrm{i}}=\mathcal{O}(1)\) when \(k_0 a=\mathcal{O}(1)\), and Fraunhofer scale
\(r_{\mathrm{F}}=2a^{2}/\lambda_0\) entering \(C_{\mathrm{f}}=\mathcal{O}(1)\)
\cite{stratton1941}. Equations
\eqref{eq:ch2.distance-regime-definitions} upgrade Eq.~\eqref{eq:ch2.distance-regime}:
they are diagnostic bands, not disjoint sharp cutoffs \cite{stratton1941}. For
\(k_0 a\ll 1\) (electrically small sources), \(C_{\mathrm{i}}\) may be taken near unity
and the intermediate zone is narrow; for extended apertures with \(k_0 a\gg 1\), the
Fresnel--Fraunhofer condition \(r\gtrsim r_{\mathrm{F}}\) controls pattern formation
\cite{harrington2001}.
Near-zone field scalings. Evaluating Eq.~\eqref{eq:ch2.outgoing-electric-volume-integral} with
Eq.~\eqref{eq:ch2.free-space-dyadic-green} for \(k_0 r\le C_{\mathrm{n}}\) yields
\begin{equation}
\label{eq:ch2.near-zone-field-scaling}
|\E(\mathbf{r})|
=
\mathcal{O}\!\big(\mu_0\omega\|\Jimp\|_{\mathcal{J}}\,k_0^{-2} r^{-3}\big)
\;+\;
\mathcal{O}\!\big(\mu_0\omega\|\Jimp\|_{\mathcal{J}}\,k_0^{-1} r^{-2}\big),
\end{equation}
with \(\|\Jimp\|_{\mathcal{J}}\) a source amplitude scale on \(\Omega_s\)
\cite{stratton1941}. Equation~\eqref{eq:ch2.near-zone-field-scaling} reflects
the \(R^{-3}\) and \(R^{-2}\) terms in Eq.~\eqref{eq:ch2.free-space-dyadic-green}; the
leading \(r^{-3}\) contribution is reactive (stored) in the sense of
Section~\ref{sec:ch01.2.4} \cite{stratton1941}. Magnetic field follows from
Eq.~\eqref{eq:ch2.outgoing-magnetic-from-curl-integral} with the same near-zone
polynomial-in-\(k_0 r\) ordering \cite{harrington2001}. On shells with
Eq.~\eqref{eq:ch1.regime-index} \(\chi(r)\gg 1\), Poynting flux through
Eq.~\eqref{eq:ch2.radiated-power-screen} is not a credible export meter
\cite{stratton1941}.
When \(C_{\mathrm{n}}<k_0 r\le C_{\mathrm{i}}(k_0 a)\), reactive and radiative terms in
Eq.~\eqref{eq:ch2.free-space-dyadic-green} are comparable:
\begin{equation}
\label{eq:ch2.intermediate-zone-balance}
|\E_{\mathrm{reactive}}(\mathbf{r})|
=
\mathcal{O}\!\big(|\E_{\mathrm{rad}}(\mathbf{r})|\,k_0 r\big),
\end{equation}
with \(\E_{\mathrm{rad}}\) the leading transverse \(e^{-\jj k_0 r}/r\) component and
\(\E_{\mathrm{reactive}}\) the remainder \cite{stratton1941}. Equation
\eqref{eq:ch2.intermediate-zone-balance} diagnoses regime mixing: storage and export coexist
in the full field pair until export is isolated by \(\mathcal{R}_\omega\)
\cite{stratton1941}. For extended sources, the Fresnel parameter
\begin{equation}
\label{eq:ch2.fresnel-parameter}
F
=
\frac{\pi a^{2}}{\lambda_0 r}
=
\frac{k_0 a^{2}}{2r}
\end{equation}
separates Fresnel (\(F\gtrsim 1\)) from Fraunhofer (\(F\ll 1\)) pattern regions on
observation shells \cite{jackson1999}. The parameter \(F\) marks the transition between
those pattern regimes; complete aperture diffraction theory is not re-developed
\cite{stratton1941}.
\begin{figure}[t]
\centering
\ChOneFigBlock{%
\begin{tikzpicture}[ch1fig, node distance=5mm and 7mm]
 \node[ch1box] (src) {Source\\$k_0 a$};
 \node[ch1box, right=7mm of src] (near) {Near\\$k_0 r\lesssim 1$};
 \node[ch1box, right=7mm of near] (mid) {Intermediate\\$F\sim 1$};
 \node[ch1box, right=7mm of mid] (far) {$1/r$\\radiative};
 \node[ch1box, below=5mm of near] (chi) {$\chi(r)\gg 1$};
 \node[ch1box, below=5mm of far] (chi2) {$\chi(r)\ll 1$};
 \draw[ch1flow] (src) -- (near) -- (mid) -- (far);
 \draw[ch1flow] (near.south) -- (near.south |- chi.north);
 \draw[ch1flow] (near.south |- chi.north) -- (chi.north);
 \draw[ch1flow] (far.south) -- (far.south |- chi2.north);
 \draw[ch1flow] (far.south |- chi2.north) -- (chi2.north);
\end{tikzpicture}%
}
\caption{Distance regimes Eq.~\eqref{eq:ch2.distance-regime-definitions}: reactive
dominance near the source, transition through Fresnel parameter
Eq.~\eqref{eq:ch2.fresnel-parameter}, radiative export in the far zone with
Eq.~\eqref{eq:ch1.regime-index} \(\chi(r)\ll 1\).}
\label{fig:ch2.distance-regime-schematic}
\end{figure}
Section~\ref{sec:ch24} \cite{stratton1941}. The far-zone reduction of integrals proceeds as follows: at distances large compared with source extent and wavelength, a pattern probe on $\Gamma_R$
samples the transverse $e^{-\jj k_0 r}/r$ wave whose angular shape is fixed by a Fourier
transform of $\Jimp$ over $\Omega_s$ \cite{jackson1999}. The volume integral at
finite $r$ still evaluates $\mathcal{S}_\omega[\Jimp]$: reactive terms propagate mathematically
but do not export watt-bearing power; only after the far-zone limit does the meter reading
coincide with $\mathcal{R}_\omega[\Jimp]$ through Eqs.~\eqref{eq:ch2.far-field-integral-reduction}--\eqref{eq:ch2.far-field-amplitude-integral}
\cite{stratton1941}. When
\(k_0 r\gg 1\) and \(r\gg r_{\mathrm{F}}\) for extended sources, outgoing dyadic
asymptotics Eq.~\eqref{eq:ch2.sommerfeld-dyadic-asymptotic} reduce
Eq.~\eqref{eq:ch2.outgoing-electric-volume-integral} to
\begin{equation}
\label{eq:ch2.far-field-integral-reduction}
\E(\mathbf{r})
=
\E_0(\theta,\phi)\,\frac{e^{-\jj k_0 r}}{r}
\;+\;\mathcal{O}(r^{-2}),
\qquad
\E_0\cdot\hat{\mathbf{r}}=0,
\end{equation}
\begin{equation}
\label{eq:ch2.far-field-amplitude-integral}
\E_0(\theta,\phi)
=
-\jj\omega\mu_0\,\frac{1}{4\pi}
\int_{\Omega_s}
\Jimp(\mathbf{r}')\,
e^{-\jj k_0\hat{\mathbf{r}}\cdot\mathbf{r}'}\,d^{3}r',
\end{equation}
with \(\hat{\mathbf{r}}=\mathbf{r}/r\) and the transversality constraint enforced by
\(\overline{\mathbf{G}}_{\infty}\) in Eq.~\eqref{eq:ch2.sommerfeld-dyadic-asymptotic}
\cite{stratton1941}. Equation
\eqref{eq:ch2.far-field-integral-reduction} is the far-zone limit of
Eq.~\eqref{eq:ch2.far-zone-transverse-radiation}; the amplitude integral
Eq.~\eqref{eq:ch2.far-field-amplitude-integral} is the vector radiation integral used in
antenna far-field post-analysis \cite{harrington2001}. The amplitude integral is what a
Fraunhofer-range pattern probe measures after $\Pi_{\mathrm{rad}}$: it Fourier-transforms
$\Jimp(\mathbf{r}')$ over the source volume while the reactive inverse-power terms that
$\mathcal{S}_\omega[\Jimp]$ carries at finite $r$ no longer enter at leading order
\cite{harrington2001}. Magnetic field in the far zone
satisfies Eq.~\eqref{eq:ch2.sommerfeld-impedance-ratio} with the same \(\E_0\)
\cite{stratton1941}.
\phantomsection\label{thm:ch2.far-field-approximation}
Let \(\Jimp\in\mathcal{J}_\Omega\) be supported in \(B_a(0)\) with outgoing radiation condition.
For \(r> r_{\min}(a,\lambda_0)\) as in Eq.~\eqref{eq:ch2.distance-regime-definitions},
the field from Eq.~\eqref{eq:ch2.outgoing-electric-volume-integral} satisfies
Eqs.~\eqref{eq:ch2.far-field-integral-reduction}--\eqref{eq:ch2.far-field-amplitude-integral}
with error \(\mathcal{O}\big((k_0 a)^{2}/(k_0 r)\,|\E_0|\,r^{-1}\big)\) for extended
sources and \(\mathcal{O}\big((k_0 r)^{-1}\,|\E_0|\,r^{-1}\big)\) for \(k_0 a\ll 1\)
\cite{stratton1941}.
Substitute Eq.~\eqref{eq:ch2.sommerfeld-dyadic-asymptotic} into
Eq.~\eqref{eq:ch2.outgoing-electric-volume-integral} and expand
\(|\mathbf{r}-\mathbf{r}'|^{-1}\) for \(|\mathbf{r}'|\le a\ll r\). The leading
\(e^{-\jj k_0 r}/r\) term factorizes into \(\E_0(\theta,\phi)\) times phase
\(e^{-\jj k_0\hat{\mathbf{r}}\cdot\mathbf{r}'}\) over the source; transversality follows
from the dyadic structure Eq.~\eqref{eq:ch2.free-space-dyadic-green} at leading order
\cite{stratton1941}. Higher orders in \(a/r\) and \(1/(k_0 r)\) give the
stated error bounds \cite{jackson1999}.
On spheres \(S_r=\{|\mathbf{r}|=r\}\), define shell power
\(P(R)=\Re\oint_{S_R}\mathbf{S}\cdot\hat{\mathbf{r}}\,dS\) as in
Eq.~\eqref{eq:ch2.radiated-power-screen} \cite{stratton1941}. In the far zone,
\begin{equation}
\label{eq:ch2.power-flux-regime-convergence}
|P(r)-P_\infty|
=
\mathcal{O}(r^{-1}),
\qquad
P_\infty
=
\frac{1}{2\eta_0}\oint_{S_r}
\big|\E_0(\theta,\phi)\big|^{2}\,d\Omega,
\end{equation} with \(\E_0\) from Eq.~\eqref{eq:ch2.far-field-amplitude-integral} and \(\eta_0\) from
Eq.~\eqref{eq:ch2.exterior-flux-asymptotic} \cite{stratton1941}. Equation
\eqref{eq:ch2.power-flux-regime-convergence} is the operator form of
Eq.~\eqref{eq:ch1.radiative-regime-condition}: \(P(r)\) stabilizes once reactive
\(r^{-2}\), \(r^{-3}\) contributions no longer dominate the shell integral
\cite{stratton1941}. Evaluating Eq.~\eqref{eq:ch2.radiated-power-screen} on
\(\mathcal{S}_\omega[\Jimp]\) before \(r\gtrsim r_c\) from
Eq.~\eqref{eq:ch1.crossover-radius} confuses storage with export
\cite{harrington2001}.
Distance regimes induce an observation-shell map on operator outputs:
\begin{equation}
\label{eq:ch2.observation-shell-regime-map}
\begin{aligned}
r\in\text{near}
&\ \Longrightarrow\
\text{meter}\ \mathcal{S}_\omega[\Jimp],\\
r\in\text{far}
&\ \Longrightarrow\
\text{meter}\ \mathcal{R}_\omega[\Jimp]\ \text{via Eqs.~\eqref{eq:ch2.far-field-integral-reduction}--\eqref{eq:ch2.far-field-amplitude-integral}},
\end{aligned}
\end{equation}
with intermediate shells requiring explicit \(\chi(r)\) from
Eq.~\eqref{eq:ch1.regime-index} before transport claims \cite{stratton1941}.
When \(k_0 a\ll 1\), Eq.~\eqref{eq:ch2.far-field-amplitude-integral} reduces to the
electric-dipole pattern proportional to
\(\widehat{\boldsymbol{\theta}}\times(\widehat{\boldsymbol{\theta}}\times\mathbf{p})\)
with \(\mathbf{p}=\int \Jimp(\mathbf{r}')\,d^{3}r'\) \cite{jackson1999}. The
near zone extends to \(r\sim\lambda_0\); far-zone reduction applies for
\(k_0 r\gg 1\) \cite{harrington2001}. When \(k_0 a\gg 1\), multiple lobes appear in
\(\E_0(\theta,\phi)\) and Fraunhofer distance \(r_{\mathrm{F}}=2a^{2}/\lambda_0\) dominates
the far-zone threshold \cite{stratton1941}. Angular-spectrum action in
Section~\ref{sec:ch29} filters $\E_0(\theta,\phi)$ without altering the integral
reduction above \cite{stratton1941}.
Equations~\eqref{eq:ch2.distance-regime-definitions}--\eqref{eq:ch2.power-flux-regime-convergence}
require linear vacuum or piecewise-constant media, fixed \(\omega\), outgoing radiation
condition, and observation outside \(\Omega_s\) \cite{stratton1941}. They fail in cutoff
waveguides where propagating modes do not exist
(Eq.~\eqref{eq:ch2.cutoff-condition}), in cavity standing-wave fields, and in lossless open
resonances without outgoing selection \cite{stratton1941}. Anisotropic beams and surface
waves require modified asymptotics beyond the locked free-space dyad template
\cite{harrington2001}. The far-field reduction asserts field asymptotics on the outgoing
branch only; it does not assert uniqueness of sources \cite{stratton1941}.
The regime partition is not a sharp geometric cutoff but a diagnostic on which terms of the same outgoing integral dominate what an instrument measures. Near and intermediate zones may appear to contradict far-zone pattern theory because substantial field energy is visible on small spheres; on the bench, that energy is largely reactive circulation diagnosed through $\mathcal{S}_\omega$, while pattern definitions attach only after $k_0 r\gg 1$ and Fresnel--Fraunhofer qualifications are met \cite{harrington2001}. Shell watt readings must not be identified with radiated power until Eq.~\eqref{eq:ch2.regime-operator-screen} is satisfied: what changes across bands is the relative weight of storage and export in the \emph{same} Maxwell solution, not the governing equations themselves \cite{poynting1884}. Subsection~\ref{sec:ch253} connects these spatial bands to flux, momentum, and operator norms imported from Chapter~\ref{ch:01}, always downstream of the integral evaluation chain established here.
